% Part: first-order-logic
% Chapter: syntax-and-semantics
% Section: formation-sequences

\documentclass[../../../include/open-logic-section]{subfiles}

\begin{document}

\olfileid{fol}{syn}{fseq}

\olsection{રચના-શ્રેણીઓ}

અનુમાનાત્મક વ્યાખ્યા વડે !!{formula}s વ્યાખ્યાયિત કરવા અને તેની
પૂરક રીત તરીકે અનુમાન વડે !!{formula}sના ગુણધર્મો સાબિત કરવા એ
સુઘડ અને કાર્યક્ષમ અભિગમ છે. જોકે !!{formula}sની રચનાને નીચેથી ઉપર,
ક્રમિક પગલાંમાં જોવી પણ ઉપયોગી થઈ શકે છે. અહીં આપણે \emph{રચના-શ્રેણી}
એ ખ્યાલ વડે એમ કરીશું.
%
પદો અને !!{formula}sને તેમની અનુમાનાત્મક વ્યાખ્યાઓનો ઉલ્લેખ કર્યા વિના
આ રીતે કેવી રીતે રજૂ કરી શકાય તે બતાવવા માટે, આપણે પહેલાં કોઈ
ભાષા~$\Lang L$માંથી લેવાયેલા સંકેતોની મનસ્વી શ્રેણીનો ખ્યાલ રજૂ કરીએ.

\begin{defn}[સંકેતશ્રેણીઓ]
\ollabel{defn:string}
ધારો કે $\Lang L$ પ્રથમ-ક્રમની ભાષા છે. \emph{$\Lang L$-સંકેતશ્રેણી}
એ~$\Lang L$ના સંકેતોની સાન્ત શ્રેણી છે. જ્યારે ભાષા~$\Lang L$ સંદર્ભથી
સ્પષ્ટ રીતે નિશ્ચિત હોય, ત્યારે આપણે ઘણી વાર $\Lang L$-સંકેતશ્રેણીને
માત્ર \emph{સંકેતશ્રેણી} કહીશું.
\end{defn}

\begin{ex}
કોઈ પણ પ્રથમ-ક્રમની ભાષા~$\Lang L$ માટે બધાં $\Lang L$-!!{formula}s
$\Lang L$-સંકેતશ્રેણીઓ છે, પરંતુ તેનું પ્રતિલોમ સાચું નથી. ઉદાહરણ તરીકે,
\[)(\Obj v_0\lif\lexists\]
એ $\Lang L$-સંકેતશ્રેણી છે, પણ $\Lang L$-!!{formula} નથી.
\end{ex}

\begin{defn}[પદો માટેની રચના-શ્રેણીઓ]
\ollabel{defn:fseq-trm}
$\Lang L$-સંકેતશ્રેણીઓની સાન્ત શ્રેણી $\tuple{t_0,\dotsc,t_n}$ને પદ~$t$
માટેની \emph{રચના-શ્રેણી} ત્યારે કહીએ જ્યારે $t \ident t_n$ હોય અને
દરેક $i \leq n$ માટે કાં તો $t_i$ !!a{variable} અથવા !!a{constant} હોય,
અથવા $\Lang L$માં કોઈ $k$-સ્થાનીય !!{function}~$f$ હોય અને એવાં
$m_1,\dotsc,m_k < i$ હોય કે
$t_i \ident f(t_{m_1},\dotsc,t_{m_k})$. જ્યારે ભેદ પાડવો જરૂરી હોય,
ત્યારે પદો માટેની રચના-શ્રેણીઓને આપણે \emph{પદ-રચના-શ્રેણીઓ} કહીશું.
\end{defn}
\sourcecorrection{OLSYN-004}{મૂળની \texttt{formation-sequences.tex}માં
$k$-સ્થાનીય વિધેય માટે $m_0,\dotsc,m_k$ અને
$f(t_{m_0},\dotsc,t_{m_k})$ લખાયેલાં છે, જે $k+1$ દલીલો આપે છે. અહીં
$k$ દલીલો સાથે સુસંગત $m_1,\dotsc,m_k$ અને
$f(t_{m_1},\dotsc,t_{m_k})$ લખાણ પુનઃસ્થાપિત કર્યું છે.}

\begin{ex}
શ્રેણી
\[
    \tuple{\Obj c_0, \Obj v_0, \Atom{\Obj f^2_0}{\Obj c_0, \Obj v_0}, \Atom{\Obj f^1_0}{\Atom{\Obj f^2_0}{\Obj c_0, \Obj v_0}}}
\]
એ પદ $\Atom{\Obj f^1_0}{\Atom{\Obj f^2_0}{\Obj c_0, \Obj v_0}}$
માટેની રચના-શ્રેણી છે; એવી જ રીતે
\[
    \tuple{\Obj v_0, \Obj c_0, \Atom{\Obj f^2_0}{\Obj c_0, \Obj v_0}, \Atom{\Obj f^1_0}{\Atom{\Obj f^2_0}{\Obj c_0, \Obj v_0}}}.
\]
પણ છે.
\end{ex}

\begin{defn}[સૂત્રો માટેની રચના-શ્રેણીઓ]
\ollabel{defn:fseq-frm}
$\Lang L$-સંકેતશ્રેણીઓની સાન્ત શ્રેણી $\tuple{!A_0,\dotsc,!A_n}$ને
$!A$ માટેની \emph{રચના-શ્રેણી} ત્યારે કહીએ જ્યારે $!A \ident !A_n$
હોય અને દરેક $i \leq n$ માટે કાં તો $!A_i$ આણ્વિક !!{formula} હોય,
અથવા એવાં $j,k < i$ અને !!a{variable}~$x$ હોય કે નીચેનામાંથી કોઈ એક
વાત સાચી હોય:
\begin{enumerate}
    \tagitem{prvNot}{$!A_i \ident \lnot !A_j$.}{}%
    \tagitem{prvAnd}{$!A_i \ident (!A_j \land !A_k)$.}{}%
    \tagitem{prvOr}{$!A_i \ident (!A_j \lor !A_k)$.}{}%
    \tagitem{prvIf}{$!A_i \ident (!A_j \lif !A_k)$.}{}%
    \tagitem{prvIff}{$!A_i \ident (!A_j \liff !A_k)$.}{}%
    \tagitem{prvAll}{$!A_i \ident \lforall[x][!A_j]$.}{}%
    \tagitem{prvEx}{$!A_i \ident \lexists[x][!A_j]$.}{}%
\end{enumerate}
જ્યારે ભેદ પાડવો જરૂરી હોય, ત્યારે સૂત્રો માટેની રચના-શ્રેણીઓને આપણે
\emph{સૂત્ર-રચના-શ્રેણીઓ} કહીશું.
\end{defn}

\begin{ex}
\[
\tuple{
    \Atom{\Obj A^1_0}{\Obj v_0},
    \Atom{\Obj A^1_1}{\Obj c_1},
    (\Atom{\Obj A^1_1}{\Obj c_1} \land \Atom{\Obj A^1_0}{\Obj v_0}),
    \lexists[\Obj v_0][(\Atom{\Obj A^1_1}{\Obj c_1} \land \Atom{\Obj A^1_0}{\Obj v_0})]
}
\]
એ $\lexists[\Obj v_0][(\Atom{\Obj A^1_1}{\Obj c_1} \land \Atom{\Obj
A^1_0}{\Obj v_0})]$ માટેની રચના-શ્રેણી છે; એવી જ રીતે
\begin{multline*}
\tuple{
    \Atom{\Obj A^1_0}{\Obj v_0},
    \Atom{\Obj A^1_1}{\Obj c_1},
    (\Atom{\Obj A^1_1}{\Obj c_1} \land \Atom{\Obj A^1_0}{\Obj v_0}),
    \Atom{\Obj A^1_1}{\Obj c_1},\\
    \lforall[\Obj v_1][\Atom{\Obj A^1_0}{\Obj v_0}],
    \lexists[\Obj v_0][(\Atom{\Obj A^1_1}{\Obj c_1} \land \Atom{\Obj A^1_0}{\Obj v_0})]
}.
\end{multline*}
પણ છે.
%
બીજા ઉદાહરણ પરથી દેખાય છે તેમ, રચના-શ્રેણીમાં ``કચરો'' હોઈ શકે છે:
એવાં !!{formula}s જે બિનજરૂરી હોય અથવા રચનામાં કોઈ ફાળો આપતાં ન હોય.
\end{ex}

\begin{prop}\ollabel{prop:formed}
$\Frm[L]$ના દરેક !!{formula}~$!A$ માટે કોઈ રચના-શ્રેણી હોય છે.
\end{prop}

\begin{proof}
ધારો કે $!A$ આણ્વિક છે. ત્યારે શ્રેણી~$\tuple{!A}$ એ $!A$ માટેની
રચના-શ્રેણી છે.
%
હવે ધારો કે $!B$ અને~$!C$ માટે અનુક્રમે રચના-શ્રેણીઓ
$\tuple{!B_0,\dotsc,!B_n}$ અને $\tuple{!C_0,\dotsc,!C_m}$ છે.
%
\begin{enumerate}
  \tagitem{prvNot}{જો $!A \ident \lnot !B$ હોય,
      તો $\tuple{!B_0,\dotsc,!B_n,\lnot !B_n}$
      એ $!A$ માટેની રચના-શ્રેણી છે.}{}
  \tagitem{prvAnd}{જો $!A \ident (!B \land !C)$ હોય,
      તો $\tuple{!B_0,\dotsc,!B_n,!C_0,\dotsc,!C_m,(!B_n \land !C_m)}$
      એ $!A$ માટેની રચના-શ્રેણી છે.}{}
  \tagitem{prvOr}{જો $!A \ident (!B \lor !C)$ હોય,
      તો $\tuple{!B_0,\dotsc,!B_n,!C_0,\dotsc,!C_m,(!B_n \lor !C_m)}$
      એ $!A$ માટેની રચના-શ્રેણી છે.}{}
  \tagitem{prvIf}{જો $!A \ident (!B \lif !C)$ હોય,
      તો $\tuple{!B_0,\dotsc,!B_n,!C_0,\dotsc,!C_m,(!B_n \lif !C_m)}$
      એ $!A$ માટેની રચના-શ્રેણી છે.}{}
  \tagitem{prvIff}{જો $!A \ident (!B \liff !C)$ હોય,
      તો $\tuple{!B_0,\dotsc,!B_n,!C_0,\dotsc,!C_m,(!B_n \liff !C_m)}$
      એ $!A$ માટેની રચના-શ્રેણી છે.}{}
  \tagitem{prvAll}{જો $!A \ident \lforall[x][!B]$ હોય,
      તો $\tuple{!B_0,\dotsc,!B_n,\lforall[x][!B_n]}$
      એ $!A$ માટેની રચના-શ્રેણી છે.}{}
  \tagitem{prvEx}{જો $!A \ident \lexists[x][!B]$ હોય,
      તો $\tuple{!B_0,\dotsc,!B_n,\lexists[x][!B_n]}$
      એ $!A$ માટેની રચના-શ્રેણી છે.}{}
\end{enumerate}
!!{formula}s પર અનુમાનના સિદ્ધાંત મુજબ દરેક !!{formula} માટે
રચના-શ્રેણી હોય છે.
\end{proof}

આપણે તેનું પ્રતિલોમ પણ સાબિત કરી શકીએ છીએ. આ મહત્વનું છે, કારણ કે
તેનાથી સૂત્રો વ્યાખ્યાયિત કરવાની આપણી બંને રીતો સમતુલ્ય છે: એ બંને
સરખાં પરિણામો આપે છે. તેનો અર્થ એ પણ થાય છે કે રચના-શ્રેણીઓની લંબાઈ
પર સામાન્ય અનુમાન વાપરીને આપણે સૂત્રો વિશેનાં પ્રમેયો સાબિત કરી
શકીએ છીએ.

\begin{lem}
\ollabel{lem:fseq-init}
ધારો કે $\tuple{!A_0,\dotsc,!A_n}$ એ~$!A_n$ માટેની રચના-શ્રેણી છે
અને $k \leq n$. ત્યારે $\tuple{!A_0,\dotsc,!A_k}$ એ~$!A_k$ માટેની
રચના-શ્રેણી છે.
\end{lem}

\begin{proof}
અભ્યાસપ્રશ્ન.
\end{proof}

\begin{prob}
\olref[fol][syn][fseq]{lem:fseq-init} સાબિત કરો.
\end{prob}

\begin{thm}
\ollabel{thm:fseq-frm-equiv}
$\Frm[L]$ એ એવી બધી $\Lang L$-સંકેતશ્રેણીઓ~$!A$નો ગણ છે કે
$!A$ માટે કોઈ સૂત્ર-રચના-શ્રેણી હોય.
\end{thm}

\begin{proof}
ભાષા~$\Lang L$ની એવી બધી સંકેતશ્રેણીઓનો ગણ $F$ લો જેમની
રચના-શ્રેણી હોય છે. \olref[fol][syn][fseq]{prop:formed}માં આપણે
$\Frm[L] \subseteq F$ જોયું છે; તેથી હવે આપણે તેનું પ્રતિલોમ
સાબિત કરીએ છીએ.

ધારો કે $!A$ની રચના-શ્રેણી $\tuple{!A_0,\dotsc,!A_n}$ છે. $n$ પર
પ્રબળ અનુમાનથી આપણે $!A \in \Frm[L]$ સાબિત કરીએ છીએ. આપણી
અનુમાન-પૂર્વધારણા એ છે કે $m < n$ લંબાઈની રચના-શ્રેણી ધરાવતી દરેક
સંકેતશ્રેણી $\Frm[L]$માં છે. રચના-શ્રેણીની વ્યાખ્યા મુજબ કાં તો
$!A \ident !A_n$ આણ્વિક છે, અથવા એવાં $j,k < n$ હોવા જોઈએ કે
નીચેનામાંથી કોઈ એક કિસ્સો બને:
\begin{enumerate}
    \tagitem{prvNot}{$!A \ident \lnot !A_j$.}{}%
    \tagitem{prvAnd}{$!A \ident (!A_j \land !A_k)$.}{}%
    \tagitem{prvOr}{$!A \ident (!A_j \lor !A_k)$.}{}%
    \tagitem{prvIf}{$!A \ident (!A_j \lif !A_k)$.}{}%
    \tagitem{prvIff}{$!A \ident (!A_j \liff !A_k)$.}{}%
    \tagitem{prvAll}{$!A \ident \lforall[x][!A_j]$.}{}%
    \tagitem{prvEx}{$!A \ident \lexists[x][!A_j]$.}{}%
\end{enumerate}
હવે આપણે કિસ્સાવાર વિચારીએ. જો $!A$ આણ્વિક હોય, તો
$!A_n \in \Frm[L]$. તેના બદલે ધારો કે
$!A \ident (!A_j \land !A_k)$. \olref[fol][syn][fseq]{lem:fseq-init}
મુજબ $\tuple{!A_0,\dotsc,!A_j}$ અને $\tuple{!A_0,\dotsc,!A_k}$ અનુક્રમે
$!A_j$ અને~$!A_k$ માટેની રચના-શ્રેણીઓ છે. આ બંને $!A$ની
રચના-શ્રેણીની ઉચિત આરંભ-ઉપશ્રેણીઓ હોવાથી તેમની લંબાઈ~$n$થી ઓછી છે.
તેથી અનુમાન-પૂર્વધારણા મુજબ $!A_j$ અને~$!A_k$, $\Frm[L]$માં છે;
અને આમ !!a{formula}ની વ્યાખ્યા મુજબ $(!A_j \land !A_k)$ પણ તેમાં છે.
બીજા કિસ્સા સમાંતર દલીલથી મળે છે.
\end{proof}
\sourcecorrection{OLSYN-005}{મૂળના આ સાબિતીમાં સામાન્ય
$\Lang L$ અને $\Frm[L]$ વચ્ચે બે સ્થળે અચાનક $\Frm[L_0]$ આવે છે.
પ્રમેયની ભાષા અને સમગ્ર દલીલ સાથે સુસંગત $\Frm[L]$ અહીં બંને
સ્થળે પુનઃસ્થાપિત કર્યું છે.}
\sourcecorrection{OLSYN-006}{મૂળના આ સાબિતીમાં સંયોજનના કિસ્સાની
પૂર્વધારણા $!A\equiv(!A_j\land !A_k)$ લખાઈ છે. અહીં તાર્કિક
સમતુલ્યતાની નહિ પણ સંકેતશ્રેણીની અભિન્નતાની જરૂર છે; વ્યાખ્યામાં
અને આ જ સાબિતીના અગાઉના ખંડમાં વપરાયેલો $\ident$ સંકેત અહીં
પુનઃસ્થાપિત કર્યો છે.}

પદો માટેની રચના-શ્રેણીઓના ગુણધર્મો !!{formula}s માટેની
રચના-શ્રેણીઓ જેવા જ છે.

\begin{prop}
\ollabel{prop:fseq-trm-equiv}
$\Trm[L]$ એ એવી બધી $\Lang L$-સંકેતશ્રેણીઓ~$t$નો ગણ છે કે
$t$ માટે કોઈ પદ-રચના-શ્રેણી હોય.
\end{prop}

\begin{proof}
અભ્યાસપ્રશ્ન.
\end{proof}

\begin{prob}
\olref[fol][syn][fseq]{prop:fseq-trm-equiv} સાબિત કરો.
સંકેત: \olref[fol][syn][fseq]{thm:fseq-frm-equiv}ના સાબિતીમાં
વાપરેલી રીત જેવી જ રીત વાપરો.
\end{prob}

રચના-શ્રેણીઓમાં બે પ્રકારનો ``કચરો'' આવી શકે છે: પુનરાવર્તિત ઘટકો
અને સૂત્ર કે પદની રચના માટે અપ્રસ્તુત ઘટકો. ન્યૂનતમ રચના-શ્રેણીઓને
જોવાથી આપણે બંનેને દૂર કરી શકીએ છીએ.

\begin{defn}[ન્યૂનતમ રચના-શ્રેણીઓ]
\ollabel{defn:minimal-fseq}
સૂત્ર~$!A$ માટેની રચના-શ્રેણી $\tuple{!A_0, \ldots, !A_n}$ને
$!A$ માટેની \emph{ન્યૂનતમ રચના-શ્રેણી} ત્યારે કહીએ જ્યારે $!A$ માટેની
બીજી દરેક રચના-શ્રેણી~$s$ માટે $s$ની લંબાઈ $n+1$થી મોટી અથવા તેના
બરાબર હોય.

તે જ રીતે, પદ~$t$ માટેની રચના-શ્રેણી $\tuple{t_0, \ldots, t_n}$ને
$t$ માટેની \emph{ન્યૂનતમ રચના-શ્રેણી} ત્યારે કહીએ જ્યારે $t$ માટેની
બીજી દરેક રચના-શ્રેણી~$s$ માટે $s$ની લંબાઈ $n+1$થી મોટી અથવા તેના
બરાબર હોય.
\end{defn}

ધ્યાન આપો કે કોઈ સૂત્ર કે પદને એક કરતાં વધુ ન્યૂનતમ રચના-શ્રેણીઓ
હોઈ શકે છે, પરંતુ એમાં બરાબર એ જ સંકેતશ્રેણીઓ હશે.

\begin{prop}
\ollabel{prop:subformula-equivs}
નીચેનાં વિધાનો સમતુલ્ય છે:
\begin{enumerate}
    \item $!B$ એ~$!A$નું ઉપ-!!{formula} છે.
    \item $!B$, $!A$ની દરેક રચના-શ્રેણીમાં આવે છે.
    \item $!B$, $!A$ની કોઈ ન્યૂનતમ રચના-શ્રેણીમાં આવે છે.
\end{enumerate}
\end{prop}

\begin{proof}
અભ્યાસપ્રશ્ન.
\end{proof}

\begin{prob}
\olref[fol][syn][fseq]{prop:subformula-equivs} સાબિત કરો.
\end{prob}

\begin{history}
રેમન્ડ સ્મલ્યને તેમના પાઠ્યપુસ્તક \emph{First-Order Logic}
\citep{Smullyan1968}માં રચના-શ્રેણીઓ રજૂ કરી હતી. રચના-શ્રેણીઓના
વધારાના ગુણધર્મો \citet{Zuckerman1973}એ સ્થાપિત કર્યા હતા.
\end{history}

\end{document}
